\documentclass[a4paper]{article}

\usepackage[utf8]{inputenc}
\usepackage[T1]{fontenc}
\usepackage{fontawesome}
\usepackage[margin=1.5cm]{geometry}
\usepackage{xcolor}
\usepackage{fancyhdr}
\usepackage{lastpage}
\usepackage{tabularx}
\usepackage{enumitem}

\pagestyle{fancy}
\fancyhf{}
\rhead{Antonio Michele Miti}
\lhead{Selected publications}
\lfoot{}
\rfoot{\thepage \hspace{1pt} / \pageref{LastPage}}
\setlength\parindent{0pt}

\usepackage{bibentry}
\usepackage{natbib}

\usepackage[hyphens]{url}
\urlstyle{same}

\usepackage[
	hidelinks,
	colorlinks = true,
	linkcolor = blue,
	urlcolor  = blue,
	citecolor = blue,
	anchorcolor = blue
]{hyperref}

\hypersetup{breaklinks=true}

\title{SELECTED PUBLICATIONS}
\date{\today}
\author{
	\href{www.antoniomiti.it}{Antonio Michele Miti},\\
	\emph{orcid: \href{https://orcid.org/0000-0002-8829-1943}
	{0000-0002-8829-1943}}.
}

\begin{document}

\nobibliography{../data/publications.bib}
\bibliographystyle{plainnat}

\maketitle

\noindent
The following publications can also be retrieved at:
\\
\url{https://www.dropbox.com/sh/asop74adf1c0gbi/AADDEeq-8XCf3IISOMzrlbroa?dl=0}.



\begin{enumerate}
	\item \bibentry{Miti2025};
	\item \bibentry{Blacker2024};
	\item \bibentry{Miti2024};
	\item \bibentry{Miti2021c};
	\item \bibentry{Miti2020}.
\end{enumerate}

\section*{Motivation for the selection}

The publications below have been selected to highlight the part of my research
most directly related to the present position. Together, they develop a coherent
research line in \textbf{multisymplectic geometry}, from symmetries and homotopy
momentum maps to higher algebras of observables and geometric reduction.
These structures provide the geometric formulation of conservation laws,
constraints, and reduced dynamics in classical field theories, and are closely
connected with the foundations of \textbf{structure-preserving numerical methods}
and multisymplectic integration. In particular, my recent work on reduction and
higher observables provides tools for investigating how geometric information
can be retained when passing from continuous models to reduced or discretized
dynamical systems.

\end{document}